\section{Trabajos relacionados}
\label{sec:estado_arte}

El criterio con que cierra la Sección~\ref{sec:marco_evaluacion} ordena también la revisión de los trabajos más
cercanos a esta tesina, que pertenecen a tres líneas de investigación: la clasificación de lesiones de piel, los
\glspl{vlm} especializados en dermatología y la generación de reportes dermatológicos. La revisión se concentra en
trabajos publicados desde 2020, mientras que los anteriores solo se incluyen cuando constituyen hitos del área, y
se indica su fecha. Varios trabajos de 2025 y
2026 están publicados solo en arXiv o medRxiv, sin revisión de pares, y en cada caso se lo indica. De cada grupo de
trabajos se presentan el problema que comparten, los métodos con que lo abordan, sus resultados y sus
limitaciones, y los detalles de arquitectura, datos y hardware se resumen en las tablas de cada línea.

\subsection{Clasificación dermatológica}
\label{sec:estado_arte_diagnostico}

La clasificación de lesiones asigna a una imagen una entre varias categorías diagnósticas, sobre imágenes
dermatoscópicas, como las de \emph{HAM10000} y los desafíos de la \gls{isic}, o sobre fotografías clínicas. Con el
criterio de
la Sección~\ref{sec:marco_generalizacion}, de cada trabajo interesa si usa una métrica que no oculte las clases
minoritarias, si mantiene las imágenes de una misma lesión en una sola partición
(Sección~\ref{sec:marco_regularizacion}), si valida en otra fuente y si se compara con clínicos. La
Tabla~\ref{tab:marco_sota_clasificacion} resume los trabajos.

En 2017, una \gls{cnn} alcanzó un desempeño comparable al de dermatólogos certificados en dos distinciones entre
lesiones malignas y benignas, carcinoma queratinocítico contra queratosis seborreica y melanoma contra
nevo~\cite{esteva2017dermatologistlevel}.
En 2019, los clasificadores del desafío ISIC 2018, entrenados con \emph{HAM10000}, obtuvieron en su prueba un
promedio de aciertos mayor que el de 511 participantes humanos, y los tres mejores superaron también a los expertos
con más de diez años de experiencia~\cite{tschandl2019comparison}. Esos dos hitos plantearon la tarea como una
comparación con el especialista, y buena parte de los trabajos posteriores se midió dentro de los desafíos.

Los ganadores de los desafíos ISIC muestran cómo se trató el desbalance de clases. \textcite{gessert2020skin}
ganaron el desafío ISIC 2019 con un ensamble de EfficientNet, SENet154 y ResNeXt, entrenado con la entropía cruzada
ponderada por la frecuencia inversa de cada clase y con particiones que mantienen juntas las imágenes de una misma
lesión. \textcite{yao2022single} buscaron, en cambio, un modelo único y de bajo costo para conjuntos pequeños y
desbalanceados: una RegNetY con una pérdida que asigna un coeficiente a cada clase, mayor para el melanoma. Los dos
usan la sensibilidad macro como métrica principal, y \citeauthor{yao2022single} muestran el intercambio que implica
favorecer una clase: en ISIC 2018, subir el coeficiente del melanoma eleva su sensibilidad de 0.772 a 0.871 a costa
de la sensibilidad macro.

La limitación es común a ambos, porque evalúan sobre la prueba del desafío. En el trabajo
de \citeauthor{gessert2020skin}, la sensibilidad macro sin metadatos baja de 72.5\,\% en validación cruzada a
63.6\,\% en la prueba oficial, una brecha que atribuyen en parte a una clase desconocida que solo aparece en ella.
\citeauthor{yao2022single} no describen una partición de validación, de modo que la selección se hace sobre la
prueba de ISIC 2018.

Otros dos trabajos llevaron la evaluación fuera del desafío y la compararon con clínicos. \textcite{liu2020deep}
extendieron la tarea de la lesión aislada a las 26 condiciones de piel más frecuentes en atención primaria, con un
ensamble que combina fotografías clínicas y la historia del paciente, entrenado con casos de teledermatología de 17
sitios de Estados Unidos y con una partición temporal y separada por paciente. Frente a 18 clínicos, su exactitud
top-1 fue de 0.67, contra 0.63 de los dermatólogos, pero entre los tumores malignos su sensibilidad top-1 fue de
0.57, contra 0.71.

\textcite{marchetti2023prospective} validaron de forma prospectiva, con biopsia como referencia,
el ganador del desafío SIIM-ISIC 2020, un ensamble de 18 modelos que, entrenados en cinco particiones, suman 90
redes por imagen. Sobre 603 lesiones y con el umbral fijado de antemano, la sensibilidad fue de 96.8\,\% y la
especificidad de 37.4\,\%, y el acceso al puntaje del modelo mejoró el \gls{auc} de 11 dermatólogos, aunque sin
modificar la atención de ningún paciente. Los dos trabajos muestran que un promedio global puede ocultar el
desempeño en las lesiones malignas y que la sensibilidad que se busca para el melanoma tiene un costo en
especificidad. Su validación, sin embargo, sigue siendo acotada: una partición temporal, que según
\citeauthor{liu2020deep} no reemplaza a un conjunto externo, o un solo centro, sin lesiones biopsiadas de pacientes
de fototipo V o VI.

Los trabajos recientes sobre \emph{HAM10000} y conjuntos afines vuelven a la comparación entre arquitecturas, sin
un protocolo común ni expertos. \textcite{kreouzi2024learning} comparan cuatro \glspl{cnn} preentrenadas en la
distinción entre melanoma y nevo de \emph{DermNet}, y suman a la exactitud el tiempo de inferencia y el tamaño del
modelo, con la mira puesta en dispositivos de pocos recursos. \textcite{aboulmira2024skinhealthmate} integran un
ensamble de dos EfficientNet en una aplicación para dermatólogos que clasifica las siete categorías de
\emph{HAM10000}. Su resultado ilustra el límite de la exactitud señalado en la
Sección~\ref{sec:marco_eval_clasif}: informan una exactitud de 93\,\%, un F1 macro de 0.81 y, en melanoma, una
sensibilidad de 0.49, y atribuyen la diferencia a las clases con pocas muestras.

\textcite{kim2025skinsavvy2} son
los únicos que usan un transformer, un Swin, y amplían \emph{HAM10000} con imágenes de revistas médicas, del archivo
ISIC y de \emph{DermNet NZ} para atenuar el desbalance, con lo que las exactitudes por clase mejoran en casi todas
las clases, salvo en la mayoritaria. Ninguno de los tres valida en otra fuente, y sus autores señalan como límites
la tarea binaria, el desempeño en pieles oscuras y la calidad de la imagen, o la diferencia entre un conjunto
curado y las imágenes tomadas en entornos no controlados.

Dos sistemas de esta línea agregan un texto generado sobre la predicción. \textcite{khan2025ensemble}, publicado
solo en arXiv, sin revisión de pares, distinguen nevo de carcinoma basocelular con un ensamble de tres \glspl{cnn} que decide
por mayoría y marca para la revisión de un especialista los casos en que sus miembros discrepan, y un modelo de
lenguaje redacta después un reporte a partir de la predicción. En SkinSavvy2, la clase predicha es la entrada de
un modelo GPT que genera recomendaciones de cuidado~\cite{kim2025skinsavvy2}. Los dos textos se analizan en la
Sección~\ref{sec:estado_arte_reportes}.

\begin{table}[!tbp]
    \centering
    \caption{Trabajos de clasificación de lesiones de piel. PCP: médicos de atención primaria; NP: enfermeros
    especializados (\emph{nurse practitioners}).}
    \label{tab:marco_sota_clasificacion}
    \footnotesize
    \setlength{\tabcolsep}{3pt}
    \begin{tabularx}{\textwidth}{@{}>{\raggedright\arraybackslash}p{1.9cm}
        c
        >{\raggedright\arraybackslash\hspace{0pt}}X >{\raggedright\arraybackslash\hspace{0pt}}X
        >{\raggedright\arraybackslash\hspace{0pt}}p{1.8cm} >{\raggedright\arraybackslash\hspace{0pt}}X
        >{\raggedright\arraybackslash\hspace{0pt}}p{1.75cm}@{}}
        \toprule
        \textbf{Trabajo} & \textbf{Año} & \textbf{Arquitectura} & \textbf{Datos} & \textbf{Métricas} & \textbf{Resultado principal} & \textbf{Expertos} \\
        \midrule
        \textcite{liu2020deep} & 2020 & Ensamble de 6 Inception-v4 con metadatos
            & 17\,777 casos de 17 sitios; 26 condiciones
            & Exactitud, sensibilidad promedio
            & Top-1 0.67; dermatólogos 0.63, PCP 0.45, NP 0.41
            & 18 participantes; 3 dermatólogos de referencia \\
        \filasep
        \textcite{gessert2020skin} & 2020 & Ensamble de CNN con metadatos
            & ISIC 2019; 5 particiones por lesión
            & Sensibilidad macro, AUC
            & Sin metadatos: 72.5\,\% en validación cruzada, 63.6\,\% en prueba
            & No \\
        \filasep
        \textcite{yao2022single} & 2022 & RegNetY única con pérdida propia
            & ISIC 2017, 2018, 2019 y \emph{7-point}
            & Sensibilidad macro y de melanoma, AUC
            & ISIC 2018: 0.875 y 0.819; ISIC 2019: 0.645
            & No \\
        \filasep
        \textcite{marchetti2023prospective} & 2023 & Ensamble de 18 modelos
            & 603 lesiones biopsiadas, 95 melanomas
            & Sensibilidad, especificidad, AUC
            & Sensibilidad 96.8\,\%; especificidad 37.4\,\%; AUC 0.857
            & 11 dermatólogos; impacto teórico \\
        \filasep
        \textcite{kreouzi2024learning} & 2024 & Cuatro CNN preentrenadas, entre ellas DenseNet121
            & \emph{DermNet}; melanoma contra nevo
            & Exactitud, AUC, tiempo de inferencia, tamaño
            & Exactitud 92.30\,\% (DenseNet121); AUC 0.957 (ResNet50V2)
            & No \\
        \filasep
        \textcite{aboulmira2024skinhealthmate} & 2024 & Ensamble de EfficientNetB1 y B5 en una aplicación
            & \emph{HAM10000}; 7 clases
            & Exactitud, F1 macro, AUC
            & Exactitud 93\,\%; F1 macro 0.81; sensibilidad de melanoma 0.49
            & No \\
        \filasep
        \textcite{khan2025ensemble} & 2025 & Ensamble de tres CNN, más LLM
            & Nevo contra carcinoma basocelular; 80/10/10 al azar
            & Exactitud, F1, pérdida logarítmica
            & Exactitud 97.67\,\% contra 93.83\,\% (ResNet50)
            & No \\
        \filasep
        \textcite{kim2025skinsavvy2} & 2025 & Swin o EfficientNetV2, más GPT
            & \emph{HAM10000} ampliado con imágenes externas
            & Exactitud, precisión, sensibilidad, F1
            & Exactitud balanceada 0.780; F1 macro 0.793 (Swin)
            & 2 dermatólogos entre 20 evaluadores \\
        \bottomrule
    \end{tabularx}
\end{table}

La línea evolucionó en tres etapas: los desafíos ISIC concentraron el trabajo sobre el desbalance, los estudios
con clínicos llevaron la evaluación fuera del desafío y los trabajos recientes volvieron a la comparación entre
arquitecturas.

Quedan cuatro limitaciones. En primer lugar, predominan las redes
convolucionales, solo \textcite{kim2025skinsavvy2} incluyen un transformer y no hay un protocolo común entre
trabajos. En segundo lugar, solo \textcite{gessert2020skin}, por lesión, y \textcite{liu2020deep}, por paciente,
describen particiones que mantienen juntos los casos relacionados, mientras que \textcite{yao2022single}
seleccionan sobre la prueba y \textcite{khan2025ensemble} parten al azar. En tercer lugar, la evaluación se ha
realizado principalmente sobre particiones o pruebas de la misma fuente, la validación fuera de ella es temporal o
de un solo centro, y ningún trabajo se desarrolla ni se valida en población de América Latina. Por último, los
sistemas de mejor desempeño son ensambles grandes, y solo \textcite{khan2025ensemble} marcan los casos
discordantes, sin medir el efecto de esa abstención.

\FloatBarrier
\subsection{Modelos visión--lenguaje especializados en dermatología}
\label{sec:estado_arte_vlm}

Los \glspl{vlm} llegan a la dermatología por las dos vías de la Sección~\ref{sec:marco_vlm}: la alineación
contrastiva, que permite clasificar sin ajuste y recuperar casos, y la unión de un codificador visual con un modelo
de lenguaje que genera texto. A esas dos vías se suma una línea de evaluación crítica, que mide los modelos
generativos en consultas reales. Las Tablas~\ref{tab:marco_sota_vlm} y~\ref{tab:marco_sota_vlm_gen} resumen los
trabajos.

En la vía contrastiva, la evolución es la del texto con que se entrena. MONET~\cite{kim2024transparent} ajusta CLIP
con 105\,550 imágenes dermatológicas y descripciones extraídas de PubMed y de libros de texto, con el fin de anotar
conceptos clínicos, como el eritema o la pigmentación, y auditar conjuntos de datos y modelos. Sus autores aclaran
que el modelo no se propone para el diagnóstico. DermLIP~\cite{yan2025derm1miccv} reemplaza ese texto por
\emph{Derm1M}, un corpus
de alrededor de un millón de pares imagen--texto organizados según una ontología cuya base crearon cuatro
dermatólogos, y en clasificación zero-shot alcanza una exactitud promedio de 0.561 sobre cuatro conjuntos,
contra 0.441 de BiomedCLIP. MAKE~\cite{yan2025makemiccai} parte de que CLIP trunca los textos a 77 tokens y
descarta detalle clínico, de modo que divide cada descripción en subtextos, los alinea con regiones de la imagen y
los pondera según su relevancia diagnóstica. Los tres se evalúan sin ajuste sobre conjuntos públicos y se comparan
con CLIP y sus variantes médicas, y solo MONET incluye dermatólogos, que verifican la recuperación de imágenes por
concepto.

Los modelos fundacionales visuales del mismo grupo aumentan la escala y agregan estudios con clínicos.
PanDerm~\cite{yan2025vision} se preentrena sin etiquetas con más de dos millones de imágenes de 11 instituciones y
cuatro modalidades, y supera al siguiente modelo en 9 de 10 conjuntos de diagnóstico.
DermFM-Zero~\cite{yan2026dermfmzero}, disponible en arXiv y Research Square, sin revisión de pares, suma un
preentrenamiento contrastivo con 1\,014\,085 pares imagen--texto para funcionar sin ajuste por tarea, y en siete
conjuntos zero-shot supera a DermLIP, con un promedio de 0.733 contra 0.561 entre la exactitud balanceada de
cinco conjuntos y el \gls{auc} de los dos de detección de melanoma. Los dos son no generativos y son
los únicos de esta línea que miden el efecto de su ayuda sobre clínicos: con DermFM-Zero, por ejemplo, 30 médicos de
atención primaria pasan de 0.266 a 0.482 de exactitud en los tres primeros diagnósticos. Su preentrenamiento
requiere cuatro H100, y sus autores reconocen que cubren solo una parte de las afecciones dermatológicas y que sus
estudios son retrospectivos o con una evaluación de equidad incompleta.

La vía generativa empezó con asistentes biomédicos de propósito general. LLaVA-Med~\cite{li2023llavamed} ajusta
LLaVA de 7\,000 millones de parámetros con figuras de PubMed Central y con conversaciones generadas por GPT-4, y se
evalúa en preguntas visuales de radiología y patología. Med-Flamingo, PeFoMed y LLaVA-Tri tampoco se evalúan en
piel~\cite{moor2023medflamingo,he2025pefomed,xie2025medtrinity25m}. SkinVL~\cite{zeng2025mmskin} adapta LLaVA-Med a
la dermatología entrenando solo la proyección visual y adaptadores LoRA, con imágenes de libros de texto y nueve
conjuntos públicos de diagnóstico reformulados como preguntas, y en zero-shot es el único de los modelos
comparados que supera el 50\,\% de exactitud en \emph{HAM10000}. La variante SkinVL-MM, que es la que usa este
trabajo, se entrena solo con el corpus de libros. Ni LLaVA-Med ni SkinVL se evaluaron con expertos.

Los trabajos más recientes de esta vía buscan que el modelo razone antes de diagnosticar, y los dos están
publicados solo en arXiv, sin revisión de pares. SkinGPT-R1~\cite{shen2025skingptr1} mantiene congelado
Vision-R1-7B y entrena solo adaptadores, con 334\,618 muestras de razonamiento sintetizado por modelos de lenguaje y
un adaptador que incorpora el tono de piel estimado. DermoGPT~\cite{ru2026dermogpt} ajusta Qwen3-VL-8B con LoRA y
después con aprendizaje por refuerzo, con una recompensa que combina la exactitud, la cercanía dentro de la
ontología diagnóstica y la coincidencia de atributos morfológicos, para seguir el orden clínico de morfología,
razonamiento y diagnóstico. Sus resultados distinguen la evaluación dentro de la distribución de entrenamiento
(ID) de la evaluación \gls{ood}. Los dos incorporan dermatólogos, para calificar las salidas del propio modelo
(SkinGPT-R1) o para revisar los datos y resolver casos como referencia humana (DermoGPT), pero no para medir el
efecto de la ayuda sobre sus decisiones. SkinGPT-R1 reconoce una latencia que impide ejecutarlo en dispositivos de
borde, y DermoGPT, una brecha que persiste con la referencia humana.

\begin{table}[!tbp]
    \centering
    \caption{\glspl{vlm} contrastivos y modelos fundacionales en dermatología. $^\dagger$: sin revisión de pares.}
    \label{tab:marco_sota_vlm}
    \footnotesize
    \setlength{\tabcolsep}{3.5pt}
    \begin{tabularx}{\textwidth}{@{}>{\raggedright\arraybackslash}p{1.9cm}
        c
        >{\raggedright\arraybackslash\hspace{0pt}}X >{\raggedright\arraybackslash\hspace{0pt}}X
        >{\raggedright\arraybackslash\hspace{0pt}}p{2.2cm} >{\raggedright\arraybackslash\hspace{0pt}}X
        >{\raggedright\arraybackslash\hspace{0pt}}p{1.9cm}@{}}
        \toprule
        \textbf{Trabajo} & \textbf{Año} & \textbf{Arquitectura} & \textbf{Datos} & \textbf{Métricas} & \textbf{Resultado principal} & \textbf{Expertos} \\
        \midrule
        \textcite{kim2024transparent} & 2024 & CLIP ajustado de forma contrastiva
            & 105\,550 imágenes con texto de PubMed y libros
            & \gls{auc} por concepto
            & Tono de piel: 0.805 contra 0.736 (CLIP)
            & 2 dermatólogos verifican la recuperación \\
        \filasep
        \textcite{yan2025derm1miccv} & 2025 & DermLIP, tipo CLIP, cuatro codificadores
            & \emph{Derm1M}, ontología con base de 4 dermatólogos
            & Exactitud, \gls{auc}, Recall@K
            & Zero-shot 0.561 contra 0.441 (BiomedCLIP)
            & Solo la ontología \\
        \filasep
        \textcite{yan2025makemiccai} & 2025 & CLIP con subtextos, alineación local y ponderación
            & 403\,563 pares imagen--texto
            & Exactitud, \gls{auc}, Recall@K
            & Zero-shot 0.493 contra 0.442 (CLIP, mismo corpus)
            & Ninguno \\
        \filasep
        \textcite{yan2025vision} & 2025 & ViT-Large auto-supervisado con maestro CLIP; no generativo
            & Más de 2 millones de imágenes, 11 instituciones
            & \gls{auc}, F1 ponderado
            & Mejora en 9 de 10 conjuntos; \emph{HAM10000} por 4.7\,\%
            & 3 estudios de lectura \\
        \filasep
        \textcite{yan2026dermfmzero}$^\dagger$ & 2026 & ViT-Large auto-supervisado y contrastivo; no generativo
            & Más de 3 millones de imágenes; 1\,014\,085 pares imagen--texto
            & F1 ponderado, \gls{auc}, sensibilidad macro, Recall@K
            & Zero-shot 0.733 contra 0.561 (DermLIP), media de exactitud balanceada y \gls{auc}; atención primaria con ayuda, de 0.266 a 0.482
            & Tres estudios, con 30, 1\,090 y 34 clínicos \\
        \bottomrule
    \end{tabularx}
\end{table}

\begin{table}[!tbp]
    \centering
    \caption{\glspl{vlm} generativos en dermatología y su evaluación crítica. $^\dagger$: sin revisión de pares. ID: dentro de la distribución; OOD: fuera de ella.}
    \label{tab:marco_sota_vlm_gen}
    \footnotesize
    \setlength{\tabcolsep}{3.5pt}
    \begin{tabularx}{\textwidth}{@{}>{\raggedright\arraybackslash}p{1.9cm}
        c
        >{\raggedright\arraybackslash\hspace{0pt}}X >{\raggedright\arraybackslash\hspace{0pt}}X
        >{\raggedright\arraybackslash\hspace{0pt}}p{2.2cm} >{\raggedright\arraybackslash\hspace{0pt}}X
        >{\raggedright\arraybackslash\hspace{0pt}}p{1.9cm}@{}}
        \toprule
        \textbf{Trabajo} & \textbf{Año} & \textbf{Arquitectura} & \textbf{Datos} & \textbf{Métricas} & \textbf{Resultado principal} & \textbf{Expertos} \\
        \midrule
        \textcite{li2023llavamed} & 2023 & LLaVA de 7B ajustado en dos etapas
            & 600\,000 pares de PubMed Central; instrucciones de GPT-4
            & Puntaje relativo con GPT-4, exactitud en VQA
            & Supera al supervisado en \emph{VQA-RAD} y \emph{PathVQA} cerradas; sin evaluar en piel
            & Ninguno \\
        \filasep
        \textcite{zeng2025mmskin} & 2025 & LLaVA-Med de 7B con LoRA
            & 171\,600 imágenes de libros y nueve conjuntos públicos (variante completa)
            & BLEU-4, BERTScore, exactitud
            & Zero-shot \emph{Patch16} 82.34\,\% contra 68.55\,\% (LLaVA-Med)
            & Ninguno \\
        \filasep
        \textcite{shen2025skingptr1}$^\dagger$ & 2025 & Vision-R1-7B congelado con adaptadores
            & 334\,618 muestras con razonamiento sintético
            & Exactitud, puntajes de 0 a 5, peor grupo
            & 82.50\,\% en \emph{Derm12345}; media de 3.6 sobre 5
            & 5 dermatólogos, 1\,000 casos \\
        \filasep
        \textcite{ru2026dermogpt}$^\dagger$ & 2026 & Qwen3-VL-8B con LoRA y aprendizaje por refuerzo
            & 772\,675 instrucciones; 33\,999 preguntas
            & Exactitud, juez automático, equidad
            & ID 78.04\,\% y OOD 66.48\,\% (Gemini-2.5-Flash: 63.37\,\% y 61.21\,\%)
            & 2 dermatólogos revisan; referencia humana en 100 casos por tarea \\
        \filasep
        \textcite{jiang2026multimodal}$^\dagger$ & 2026 & Evaluación sin ajuste de cinco modelos
            & Tres conjuntos públicos y 5\,811 consultas de un hospital
            & Exactitud de los tres primeros diagnósticos, triaje
            & Imagen sola en consulta real: 1.50\,\% a 13.35\,\% (abiertos), 24.65\,\% (GPT-4.1)
            & 3 dermatólogos validan el juez \\
        \bottomrule
    \end{tabularx}
\end{table}

La línea de evaluación crítica mide el desempeño de estos modelos en la consulta real.
\textcite{jiang2026multimodal}, publicado solo en arXiv, sin revisión de pares, y de un grupo independiente de los
anteriores, evalúan sin ajuste cinco modelos, entre ellos LLaVA-Med, SkinGPT-4
(Sección~\ref{sec:estado_arte_reportes}) y GPT-4.1, en diagnóstico
diferencial y triaje, sobre tres conjuntos públicos y 5\,811 consultas de un hospital. Con la imagen sola, la
exactitud de los tres primeros diagnósticos en las consultas reales va de 1.50\,\% a 13.35\,\% en los modelos
abiertos y llega a 24.65\,\% en GPT-4.1. El texto del médico remitente mejora los resultados, pero los vuelve
sensibles a los errores de ese texto, y los autores concluyen que no corresponde usar estos modelos sin
supervisión.

La vía contrastiva avanzó de adaptar CLIP con texto de la literatura a corpus dermatológicos con ontología y a
preentrenamientos visuales a gran escala, y la generativa, de asistentes biomédicos que no se evaluaron en piel a
modelos adaptados a ella y a modelos que razonan antes de diagnosticar.

En esta línea se identifican cuatro
limitaciones. Ningún trabajo combina un clasificador supervisado con un generador ni recupera casos parecidos para
anclar el texto. La incertidumbre aparece como limitación, o en DermoGPT para ponderar sus predicciones, y no como
una decisión entre afirmar, ofrecer diferenciales o derivar (Sección~\ref{sec:marco_tres_salidas}). El cómputo va
de ocho RTX 4090 para ajustar SkinGPT-R1 a cuatro H100 para preentrenar PanDerm o DermFM-Zero. Por último, los
estudios con clínicos son pocos, y solo PanDerm y DermFM-Zero miden el efecto de la ayuda sobre sus decisiones.

\FloatBarrier
\subsection{Generación de reportes dermatológicos}
\label{sec:estado_arte_reportes}

Un reporte generado se juzga por los dos requisitos de la Sección~\ref{sec:marco_rep_ejes}, y los trabajos de esta
línea difieren en cómo lo producen, de extremo a extremo o condicionado en la salida de otro modelo
(Sección~\ref{sec:marco_rep_generacion}), en la evidencia que usan y en cómo evalúan el texto. La
Tabla~\ref{tab:marco_generadores} resume los generadores, y la Tabla~\ref{tab:marco_sota_reportes}, los trabajos
que se ocupan de la evaluación y de la recuperación.

Los primeros generadores dermatológicos de extremo a extremo buscan dar acceso a una consulta sin especialista.
SkinGPT-4~\cite{zhou2024pretrained} ajusta MiniGPT-4 con Llama-2-13B en dos pasos, primero para alinear la imagen
con conceptos clínicos y después para aprender diagnósticos a partir de notas de médicos, de modo que cualquier
persona cargue una fotografía y obtenga en una conversación una descripción de la lesión, un diagnóstico y
recomendaciones.
DermatoLlama~\cite{yilmaz2025resourceefficient}, publicado solo en medRxiv, sin revisión de pares, persigue en
cambio que el modelo se ejecute localmente, por privacidad y por costo. Ajusta con LoRA un Llama~3.2 de 11 mil
millones de parámetros con 1.2 millones de textos sintéticos (\emph{DermaSynth}), que modelos más grandes redactaron
sobre imágenes abiertas, y se ejecuta en una sola GPU.

Los dos trabajos difieren sobre todo en la evaluación. En
SkinGPT-4, cinco dermatólogos certificados juzgan 150 casos reales, sin métricas con referencia, y consideran
correcto o pertinente el 80.63\,\% de los diagnósticos. DermatoLlama, en cambio, mide la exactitud del diagnóstico
y el parecido con referencias anotadas por tres dermatólogos coautores, y ningún experto revisa los textos
generados. Entre sus limitaciones, los autores de SkinGPT-4 reconocen la poca representación de los fototipos V y VI
y el riesgo de que las alucinaciones den información errónea al paciente, y los de DermatoLlama, que los datos
sintéticos pueden perder rasgos sutiles en casos raros. SkinVL y SkinGPT-R1, descriptos en la
Sección~\ref{sec:estado_arte_vlm}, pertenecen también a este grupo.

Skin-R1~\cite{liu2025skinr1}, publicado solo en arXiv, sin revisión de pares, se ocupa de que el diagnóstico venga
con un razonamiento fundado y no de redactar un reporte completo, el mismo objetivo de SkinGPT-R1 y DermoGPT. Sobre
Qwen2.5-VL-7B con LoRA, genera trayectorias de razonamiento a partir de libros de texto, con la jerarquía de
enfermedades y el diagnóstico diferencial, y las usa para un ajuste supervisado seguido de aprendizaje por refuerzo.
Se evalúa con preguntas de opción múltiple, un formato en el que SkinVL casi nunca produce respuestas válidas, y
alcanza una exactitud media de 0.7171 en un conjunto externo de dermatoscopía. No mide la calidad del texto del
razonamiento ni incluye expertos.

Un segundo grupo separa el diagnóstico de la redacción, como la generación condicionada de la
Sección~\ref{sec:marco_rep_generacion}: un clasificador decide y un modelo de lenguaje escribe. Los dos sistemas de
este grupo, presentados en la Sección~\ref{sec:estado_arte_diagnostico}, redactan con modelos grandes a partir de la
predicción y no de la imagen. En el de \textcite{khan2025ensemble}, Llama-3 de 70 mil millones de parámetros,
consultado por API, recibe el nombre de la enfermedad, el tipo de consenso entre los clasificadores y la confianza, y
redacta la descripción de la lesión, los síntomas por vigilar, el tratamiento y las señales de alarma, de modo que la
descripción no proviene de la imagen. El texto solo se valora con una rúbrica propia de contenido. En
SkinSavvy2~\cite{kim2025skinsavvy2}, GPT-4 combina la clase predicha con la edad, el sexo, el tipo de piel y otros
datos del usuario para redactar una recomendación de cuidado, que puntuaron 20 participantes, dos de ellos
dermatólogos. Ninguno de los dos compara sus textos con reportes clínicos de referencia.

\textcite{chakrabarti2025skin} abordan la falta de pares imagen--reporte mediante destilación. GPT-4V actúa como
maestro: a partir de la imagen, una descripción de libro de texto, los rasgos anotados de la lista de 7 puntos y el
diagnóstico, produce pseudorreportes con hallazgos, impresión y diagnóstico diferencial, con los que se ajusta
LLaVA-1.5-13B sobre 305 imágenes dermatoscópicas. Se compara también con dos variantes que anteponen al
prompt los pseudorreportes más parecidos, recuperados con BioMedCLIP o PubMedCLIP, y que obtienen puntajes
menores que el modelo ajustado. Evalúa además los rasgos de la lista y la sensibilidad y especificidad para
melanoma, y un médico revisa 30 reportes. Como la referencia la escribe el propio maestro, sin embargo, las métricas
con referencia miden el parecido con GPT-4V. Los autores atribuyen a la estructura de plantilla de los
pseudorreportes un SacreBLEU inusualmente alto y señalan que las métricas léxicas no bastan.

\begin{table}[!tbp]
    \centering
    \caption{Generadores de texto dermatológico. «Métricas con referencia»: BLEU, ROUGE y BERTScore; «---»: no
    informado.}
    \label{tab:marco_generadores}
    \small
    \begin{tabularx}{\textwidth}{@{}>{\raggedright\arraybackslash}p{3.2cm}
        >{\raggedright\arraybackslash\hspace{0pt}}X >{\raggedright\arraybackslash\hspace{0pt}}X l@{}}
        \toprule
        \textbf{Trabajo} & \textbf{Modelo que redacta} & \textbf{Evaluación del texto} & \textbf{Hardware} \\
        \midrule
        SkinGPT-4~\cite{zhou2024pretrained} & MiniGPT-4 con Llama-2-13B & dermatólogos sobre 150 casos reales & 8 A100 \\
        \filasep
        SkinVL~\cite{zeng2025mmskin} & LLaVA-Med con LoRA & métricas con referencia & 8 RTX 3090 \\
        \filasep
        DermatoLlama~\cite{yilmaz2025resourceefficient} & Llama 3.2 11B con LoRA & exactitud y métricas con referencia & 4 RTX 4090 \\
        \filasep
        Skin-R1~\cite{liu2025skinr1} & Qwen2.5-VL-7B & exactitud del diagnóstico en opción múltiple & 2 A100 \\
        \filasep
        SkinGPT-R1~\cite{shen2025skingptr1} & Vision-R1-7B congelado con adaptadores & puntajes de un juez automático y de 5 dermatólogos & 8 RTX 4090 \\
        \filasep
        Khan et al.~\cite{khan2025ensemble} & Llama-3 70B por API & --- & --- \\
        \filasep
        SkinSavvy2~\cite{kim2025skinsavvy2} & GPT-4 & encuesta a 20 participantes, 2 dermatólogos & 4 RTX A6000 \\
        \filasep
        Chakrabarti et al.~\cite{chakrabarti2025skin} & LLaVA-1.5-13B destilado de GPT-4V & métricas con referencia, corrección diagnóstica y revisión de un médico & --- \\
        \bottomrule
    \end{tabularx}
\end{table}

\begin{table}[!tbp]
    \centering
    \caption{Evaluación de narrativas dermatológicas y recuperación de contexto para \glspl{vlm} médicos.}
    \label{tab:marco_sota_reportes}
    \footnotesize
    \setlength{\tabcolsep}{3.5pt}
    \begin{tabularx}{\textwidth}{@{}>{\raggedright\arraybackslash}p{1.9cm}
        c
        >{\raggedright\arraybackslash\hspace{0pt}}X >{\raggedright\arraybackslash\hspace{0pt}}X
        >{\raggedright\arraybackslash\hspace{0pt}}p{2.2cm} >{\raggedright\arraybackslash\hspace{0pt}}X
        >{\raggedright\arraybackslash\hspace{0pt}}p{1.9cm}@{}}
        \toprule
        \textbf{Trabajo} & \textbf{Año} & \textbf{Arquitectura} & \textbf{Datos} & \textbf{Métricas} & \textbf{Resultado principal} & \textbf{Expertos} \\
        \midrule
        \textcite{shen2025dermbench} & 2025 & Juez automático y evaluador LLaVA v1.5 7B con LoRA
            & 4\,000 imágenes de \emph{DermNet} con referencia; 500 imágenes y nueve modelos
            & Error absoluto medio contra médicos, seis dimensiones
            & Error medio 0.363 (DermBench) y 0.178 (DermEval)
            & 2 dermatólogos, coautores \\
        \filasep
        \textcite{xia2025mmedrag} & 2025 & LLaVA-Med-1.5 con recuperación de reportes y optimización de preferencias
            & Radiología, oftalmología y patología
            & Exactitud, F1, \gls{auc}; BLEU, ROUGE-L, METEOR
            & \emph{IU-Xray}: exactitud 75.47\,\% a 89.54\,\%; BLEU 9.64 a 31.38
            & No \\
        \filasep
        \textcite{xia2024rule} & 2024 & LLaVA-Med-1.5 con recuperación de reportes, $k$ calibrado y optimización
            de preferencias
            & Radiología y oftalmología
            & Exactitud, F1; BLEU, ROUGE-L, METEOR
            & \emph{IU-Xray}: exactitud 75.47\,\% a 87.84\,\%; BLEU 9.64 a 27.53
            & No \\
        \bottomrule
    \end{tabularx}
\end{table}

Dos líneas complementan a los generadores. La primera se ocupa de cómo evaluar el texto.
\textcite{shen2025dermbench}, presentado en ECCV 2026, parten de que las métricas de similitud textual, según los
autores, correlacionan poco con la utilidad clínica, y proponen dos evaluadores. DermBench es un juez automático que
no ve la imagen y compara el texto con una narrativa de referencia redactada por modelos de lenguaje a partir de la
etiqueta verdadera y corregida por dermatólogos. DermEval es un evaluador sin referencia, ajustado con aprendizaje
por refuerzo contra los puntajes de médicos. Sobre nueve modelos, el error absoluto medio respecto de los médicos es
0.363 para DermBench y 0.178 para DermEval. Los puntajes de referencia los asignaron dos dermatólogos coautores, sin
que se cuantificara el acuerdo entre ellos, y los autores reconocen que las referencias redactadas con modelos pueden
introducir sesgo.

La segunda línea condiciona la generación con contexto recuperado. MMed-RAG~\cite{xia2025mmedrag} y
RULE~\cite{xia2024rule}, del mismo grupo, buscan mejorar la factualidad de LLaVA-Med-1.5 con reportes de texto
recuperados a partir de la imagen, en radiología, oftalmología y, en MMed-RAG, patología. Parten del mismo problema:
el modelo puede ignorar la imagen, o un contexto equivocado puede degradar respuestas que sin él eran correctas. RULE
calibra cuántos reportes recuperar con un control estadístico del riesgo, y MMed-RAG corta la lista donde la
similitud cae de forma brusca; los dos ajustan además el generador con optimización directa de preferencias, para que
use el contexto solo cuando ayuda. En \emph{IU-Xray}, la exactitud de las preguntas sube de 75.47\,\% a 87.84\,\%
con RULE y a 89.54\,\% con MMed-RAG, pero en la ablación de RULE el contexto recuperado, agregado sin más, empeora
\emph{MIMIC-CXR} de 75.79\,\% a 67.35\,\%. Ninguno de los dos recupera imágenes, prevé la abstención, incluye expertos ni
trata la dermatología.

En esta línea, el texto se evalúa de tres maneras que casi no se combinan (Sección~\ref{sec:marco_eval_texto}). Las
métricas con referencia miden el parecido con un texto y no la corrección clínica, y las referencias suelen ser
sintéticas o anotadas por los propios autores. La corrección del diagnóstico (Skin-R1, DermatoLlama) deja fuera la
descripción. Los expertos son pocos: cinco dermatólogos en SkinGPT-4, dos entre 20 participantes en SkinSavvy2, un
médico en el trabajo de \citeauthor{chakrabarti2025skin} y dos coautores en DermBench. La recuperación puede mejorar
la factualidad y también degradarla, y en todos los casos se recuperan textos y no imágenes. Por último, varios
sistemas dependen de modelos cerrados o consultados por API para redactar o para generar sus datos o referencias, y
el hardware llega a ocho A100.

\FloatBarrier
\subsection{Posicionamiento de la propuesta}
\label{sec:estado_arte_posicion}

La revisión de las tres líneas permite identificar limitaciones que aún no han sido abordadas de manera conjunta.
La evaluación de los clasificadores de lesiones se ha realizado principalmente sobre la prueba de un desafío, un
solo centro o una partición temporal, y los de mejor desempeño son ensambles grandes. Entre los \glspl{vlm}, los
contrastivos especializados clasifican y recuperan casos pero no generan texto, los generativos requieren un
cómputo elevado y los \glspl{vlm} médicos de propósito general obtienen en piel exactitudes menores que las de los
modelos adaptados a ella. En la generación de reportes, la mayoría de los sistemas se evalúa por parecido textual o
por la exactitud del diagnóstico, y la recuperación estudiada trae textos de otras especialidades.

A partir de las limitaciones identificadas, se propone un \emph{pipeline} que integra un clasificador que estima el
diagnóstico, una recuperación de casos parecidos por imagen y un generador de reportes, con una regla de decisión
que fija cuánto afirma el texto (Sección~\ref{sec:marco_tres_salidas}). La Tabla~\ref{tab:marco_brechas} relaciona
cada limitación con la decisión del trabajo que la atiende. Como advierten MMed-RAG y RULE, un contexto equivocado
puede degradar la respuesta~\cite{xia2025mmedrag,xia2024rule}, y por eso el análisis incluye cuánto degrada el
reporte un vecino equivocado.

El codificador contrastivo es DermLIP. DermFM-Zero lo supera en su propio protocolo, pero DermLIP, en su variante
ViT-B/16, se ejecuta en una GPU de consumo, y un mismo espacio sirve para clasificar sin ajuste, entrenar un
clasificador lineal y recuperar vecinos visuales (Sección~\ref{sec:marco_recuperacion_imagenes},
\cite{yan2025derm1miccv}). Por la misma restricción de cómputo, los modelos generativos se adaptan con LoRA
(Sección~\ref{sec:marco_llm_lora}).

Dos trabajos cubren una parte de esta propuesta. El ensamble de \textcite{khan2025ensemble} deriva a un
especialista los casos en que sus miembros no coinciden, pero la concordancia solo agrega al texto un aviso de
revisión, y el texto lo escribe un modelo consultado por API que no ve la imagen. \textcite{chakrabarti2025skin}
miden la corrección diagnóstica de los reportes, pero de un único modelo destilado. Ninguno combina en una cascada
la recuperación de casos, la decisión por concordancia y la generación del reporte. La implementación y el
protocolo de evaluación se describen en el Capítulo~\ref{cp:metodos}.

\begin{table}[H]
    \centering
    \caption{De las brechas de la revisión a las decisiones del trabajo.}
    \label{tab:marco_brechas}
    \small
    \begin{tabularx}{\textwidth}{@{}>{\raggedright\arraybackslash}p{3.3cm}
        >{\raggedright\arraybackslash\hspace{0pt}}X >{\raggedright\arraybackslash\hspace{0pt}}X@{}}
        \toprule
        \textbf{Brecha} & \textbf{Evidencia en la revisión} & \textbf{Decisión del trabajo} \\
        \midrule
        Sin comparación bajo un mismo protocolo
            & Clasificadores casi solo convolucionales, validados sobre pruebas de desafíos o en un centro
              (Sección~\ref{sec:estado_arte_diagnostico}); \glspl{vlm} medidos entre sí o con clínicos
              (Sección~\ref{sec:estado_arte_vlm})
            & Comparar clasificadores supervisados y modelos médicos generales y dermatológicos con un mismo
              protocolo, con partición disjunta por lesión y un conjunto externo \\
        \filasep
        Texto evaluado por su parecido con una referencia
            & Métricas con referencia, exactitud del diagnóstico o un juez que no ve la imagen
              (Sección~\ref{sec:estado_arte_reportes})
            & Medir por separado la fidelidad diagnóstica y la calidad descriptiva
              (Sección~\ref{sec:marco_rep_ejes}) \\
        \filasep
        Incertidumbre y evidencia sin integrar
            & La incertidumbre como limitación y no como decisión; contexto recuperado de otras especialidades
              (Secciones~\ref{sec:estado_arte_vlm} y~\ref{sec:estado_arte_reportes})
            & Recuperar casos parecidos por imagen, usar la concordancia como señal de confianza y decidir entre
              afirmar, ofrecer diferenciales o derivar
              (Secciones~\ref{sec:marco_integracion} y~\ref{sec:marco_tres_salidas}) \\
        \filasep
        Cómputo fuera del alcance de un equipo sin infraestructura
            & Ocho RTX 4090, cuatro H100, ensambles de 90 redes o un modelo de 70 mil millones de parámetros
              por API
            & Ejecutar todo en una GPU de consumo, con un codificador contrastivo dermatológico y LoRA \\
        \bottomrule
    \end{tabularx}
\end{table}

\FloatBarrier
